\documentclass[11pt]{article}

% Glossary of formal terminology used across the Hart-Forcade,
% Hart-Nelson-Forcade, Morgan-Hart-Forcade, and Shinohara-Seko papers,
% and the Fortran enumlib codebase. Populated during Phase 4 paper-reading; this
% file currently contains the structure and a few seed entries to be
% expanded.

\usepackage{amsmath, amssymb}
\usepackage{geometry}
\geometry{margin=1in}
\usepackage{enumitem}
\usepackage{hyperref}

\title{Enumlib.jl --- Terminology Glossary}
\author{Built from Hart, Forcade, Nelson, Morgan, Shinohara, and Seko papers}
\date{(Living document. Phase 4 of the pre-refactor research.)}

\begin{document}
\maketitle

\section*{Purpose}

Terms used across the enumeration literature have specific meanings that
are worth pinning down precisely before any refactor is designed. This
glossary is the single source of truth for that vocabulary in the
Enumlib.jl repo --- when a design doc or a function name uses one of
these terms, it means \emph{this} thing.

Entries below are stubs to be filled in during Phase 4 (paper digests).
Each entry should cite the paper(s) and Fortran/Julia routine(s) where
the term originates and is used.

\section{Core combinatorial concepts}

\begin{description}[leftmargin=2em, style=nextline]
\item[Derivative structure / derivative superstructure]
A supercell of a parent (Bravais) lattice together with a periodic
decoration of its atomic sites. ``Derivative'' because the supercell
basis vectors are integer linear combinations of the parent's basis
vectors. (Hart \& Forcade 2008, Sec.~I.)

\item[Parent lattice $L$]
The reference Bravais lattice, treated as an infinite group under
addition. Sites in the parent are points of $L$ (plus, for multilattices,
displacements by basis vectors --- see ``dset'' below).

\item[Superlattice / supercell $S$]
A sublattice of $L$ of finite \emph{index} $n$ --- i.e., $L/S$ has $n$
cosets. Geometrically: a supercell containing $n$ parent cells.

\item[Index $n$]
$|L/S|$. Equals $\det(H)$ for the supercell HNF $H$, which equals the
volume ratio of supercell to parent cell.

\item[HNF (Hermite Normal Form)]
The lower-triangular integer-matrix representative of a superlattice's
basis-change matrix. In 3D:
\[
H = \begin{pmatrix} a & 0 & 0 \\ b & c & 0 \\ d & e & f \end{pmatrix},
\quad 0 \le b < c, \quad 0 \le d, e < f.
\]
$\det H = acf = n$. Each superlattice (as a set) has a unique HNF.

\item[SNF (Smith Normal Form)]
The diagonal integer-matrix decomposition $S = LHR$ with $L, R$
unimodular and $S = \mathrm{diag}(s_1, s_2, s_3)$ with $s_1 \mid s_2 \mid s_3$.
Provides the quotient group $L/S$ as a direct sum
$\mathbb{Z}_{s_1} \oplus \mathbb{Z}_{s_2} \oplus \mathbb{Z}_{s_3}$.
\textbf{Caution:} SNF is unique but the transformation matrices $L, R$
are not. This non-uniqueness drove the
\texttt{compare\_two\_enum\_files} machinery in Fortran enumlib v2.0.3.

\item[Quotient group $G = L/S$]
Finite group of order $n$ formed by the cosets of $S$ in $L$. The set on
which labelings live; all duplicate-elimination in Hart-Forcade 2008
happens inside $G$.

\item[Labeling / coloring]
A function $\sigma : G \to \{0, 1, \ldots, k-1\}$. Equivalently, an
$n$-digit base-$k$ integer when $G$'s elements are linearly indexed.
``Coloring'' tends to be used in the Julia code; ``labeling'' in the
Fortran code; same concept.

\item[Multilattice]
A parent unit cell with more than one basis atom (e.g., HCP has 2,
perovskite has 5). The site count per supercell is $n \cdot n_D$ where
$n_D$ is the number of basis atoms.
(See Hart \& Forcade 2009 for the multilattice treatment.)

\item[dset (d-vector set)]
The set $\{d_1, \ldots, d_{n_D}\}$ of basis-atom displacements within a
parent cell. Symmetry operations of the parent lattice act on the dset
as permutations. (See \texttt{multilattice\_dset\_mapping\_writeup.pdf}.)
\end{description}

\section{Group-theoretic concepts}

\begin{description}[leftmargin=2em, style=nextline]
\item[Point group / lattice point group]
The set of orthogonal transformations (rotations + reflections) that
map the lattice to itself. For fcc/bcc/sc: 48 elements. For hex: 24.
For tetragonal: 16. Lower symmetry $\Rightarrow$ more inequivalent
HNFs at a given index.

\item[Permutation group of a supercell $H$]
The set of permutations of supercell sites induced by symmetry
operations of the parent lattice that fix $H$. Built from the
\texttt{RotPermList.perm} array in Fortran, computed by
\texttt{getPermG} in Julia. The action on labelings is what drives
duplicate elimination.

\item[Translation group]
The cyclic subgroup of the supercell's permutation group corresponding
to translations (vs.\ rotations). One element per coset of $L'$ in $L$,
so $|T| = n$. Always a subgroup of the full permutation group; in the
Fortran code, sorted to come first via \texttt{getTransGroup}.

\item[Stabilizer subgroup of a coloring (or partial coloring)]
The subgroup of the symmetry group that leaves a (possibly partial)
coloring invariant. Used to test uniqueness: a coloring is symmetry-
equivalent to one already visited iff some stabilizer element maps
its location vector to a smaller one (lex order). Shrinks
monotonically with tree depth in the Morgan-Hart-Forcade 2017
algorithm, which is what gives that algorithm its speed.

\item[Partial coloring]
An intermediate node in the enumeration tree (Morgan-Hart-Forcade
2017) where only the first $\ell$ of $k$ colors have been placed.
All full colorings beneath a partial coloring share the
already-placed colors, so a single equivalence check on the partial
coloring eliminates an entire subtree.

\item[Recursive stabilizer (algorithm)]
The Morgan-Hart-Forcade 2017 enumeration approach: traverse a tree
of partial colorings, prune at each node by checking the location
vector under the parent's stabilizer subgroup, recompute the
stabilizer for each new unique partial coloring, recurse. Sub-linear
in the number of full configurations; the alternative to crossing-out.

\item[g-space / group-space coordinates]
Coordinates of a supercell site as an element of the quotient group
$G = L/L'$. Computed via the homomorphism $h(x) = [U A^{-1} x]_S$
from the SNF, with $S$-modular reduction. The unifying coordinate
system on which all the duplicate-elimination machinery operates.

\item[Pólya enumeration (Pólya--Redfield theorem)]
A counting principle: for a finite group $H$ acting on a set $X$, the
number of inequivalent $k$-colorings is $\frac{1}{|H|} \sum_{\rho \in H}
k^{c(\rho)}$ where $c(\rho)$ is the cycle count of $\rho$ on $X$. The
fixed-concentration extension is in Hart-Nelson-Forcade 2012,
Appendix A.2.

\item[Burnside averaging]
The $\frac{1}{|H|}\sum_{\rho \in H}$ part of the Pólya formula --- counting
fixed colorings per group element and averaging. Sometimes used
interchangeably with ``Pólya enumeration'' even though they differ in
how they count fixed configurations.
\end{description}

\section{Equivalence and reduction}

\begin{description}[leftmargin=2em, style=nextline]
\item[Symmetry-equivalent / inequivalent]
Two derivative structures are equivalent if they're related by a
combination of (a) a parent-lattice rotation, (b) a translation, (c) a
relabeling (e.g., $a \leftrightarrow b$). Enumerating ``unique''
derivative structures means listing one representative per equivalence
class.

\item[Translation duplicate]
Two labelings that differ only by adding some $t \in G$ to every site.
The $n$ translations partition the $k^n$ labelings into orbits of size $n$
generically. Eliminated in step 5b of Hart-Forcade 2008.

\item[Label-exchange duplicate]
Two labelings related by a relabeling permutation on $\{0, \ldots, k-1\}$
(e.g., swapping $a \leftrightarrow b$). For binary, halves the labeling
count. Step 5c.

\item[Super-periodic labeling]
A labeling fixed by some non-identity translation in $G$ --- equivalently,
its actual period is shorter than the supercell. Such a structure is
already enumerated as a smaller-index derivative; eliminate. Step 5d.

\item[Label-rotation duplicate]
Two labelings of the \emph{same} HNF related by a parent-lattice rotation
$R$ that fixes the superlattice. Eq.~3 of Hart-Forcade 2008 expresses the
induced permutation on $G$:
\[
G' = L A^{-1} R (L A^{-1})^{-1} G
\]
entirely inside the quotient group, no geometric reference needed.
Step 6.

\item[Complementary labeling]
A labeling and its label-flipped counterpart (e.g., $aabb \leftrightarrow
bbaa$). Special case of label-exchange. Relevant in
\texttt{compare\_two\_enum\_files} when comparing different-revision
outputs.

\item[Multipermutation]
Forcade's combinatorial trick for enumerating distinct ordered
arrangements of multisets in $O(N)$ time. Underlies the labeling
enumeration in fixed-concentration cases.
(See \texttt{multiperms.pdf}; original work in Forcade 2010 notes.)

\item[Inactive site]
A dset position whose set of allowed species has size 1 --- the species
is locked, contributing no configurational freedom. Inactive sites are
stripped from the labeling space before enumeration starts and re-inserted
into the output afterward with their hardcoded species. Examples: the
oxygen sites in a perovskite \mbox{ABO\textsubscript{3}} enumeration, where
A and B sublattices are substitutional but the O sublattice is fixed.
The convention originates with the UNCLE code; the same name is kept
here.

\item[Equivalent sites]
Two distinct dset positions that the user declares must always carry the
\emph{same} species label across configurations, even though both
sites individually retain configurational freedom. Both sites
participate in the enumeration but the dispatcher passes only one
representative of each equivalence class to the algorithm; after
enumeration, every member of the class receives the canonical's label.
This is distinct from inactive sites: equivalent sites \emph{have} a
choice (multiple allowed species), but they share whatever choice is
made.

\item[Equivalencies]
The user-declared partition of dset positions into equivalence classes
of equivalent sites; the union of all \texttt{equate!(sites, i, j)}
calls (or, equivalently, the upfront \texttt{Sites(list, classes)}
constructor argument). Equivalencies are \emph{not} derived from the
parent space group --- they capture user-asserted constraints that the
parent symmetry doesn't see, such as mirror-image layers in a slab
geometry where the vacuum has broken 3D periodicity but the user knows
the layers should share their composition.
\end{description}

\section{Concentration / site restriction}

\begin{description}[leftmargin=2em, style=nextline]
\item[Concentration / multiplicity vector $(a_1, \ldots, a_k)$]
The number of sites of each species, $\sum_i a_i = n$. Specifies a single
target stoichiometry.

\item[Concentration range]
Per-species triple $[\text{min\_num}, \text{max\_num}, \text{denom}]$
specifying allowed fractions $[\text{min}/\text{denom}, \text{max}/\text{denom}]$.
A single range expands into a list of integer multiplicity vectors at
each cell volume.

\item[Multinomial coefficient $C$]
$C = \binom{n}{a_1, a_2, \ldots, a_k} = n! / \prod_i a_i!$. The number
of distinct permutations of a multiset with multiplicities $(a_1, \ldots, a_k)$.
For binary at the median concentration, $C \sim 2^n / \sqrt{n}$.

\item[Multiset / mixed-radix hash]
Hart-Nelson-Forcade 2012 perfect hash from a fixed-concentration
configuration to a consecutive integer in $[0, C-1]$. Per-color rank
$x_i$ + mixed-radix combine: $y = x_1 + C_1 (x_2 + C_2 (\ldots))$.
Crucial enabling step for fixed-concentration enumeration at large $n$.

\item[Concentration-restricted enumeration]
Enumerating only the configurations whose multiplicities fall in a
specified concentration range. Avoids the $k^n$ memory wall when
concentrations are narrow.

\item[Site restriction]
A per-site constraint excluding certain species (e.g., site 4 cannot
be red). Combined with concentration restrictions, switches the
algorithm from indexed hash enumeration to backtracking tree search
with branch pruning. Encoded as a mask matrix.

\item[Backtracking tree search]
The site-restricted-plus-concentration-restricted enumeration
strategy. Each tree level fixes the color at one site; branches
violating site or concentration constraints are pruned, skipping
entire sub-branches.

\item[Inactive / spectator site]
A site with only one allowed species --- no configurational freedom.
Stripped from the enumeration on entry, re-inserted before output.
Introduced in Fortran enumlib v2.0.0. Conceptually a degenerate
limit of site restriction (only one species allowed).

\item[Equivalent site]
A site declared identical to a canonical site (e.g., images on the
two sides of a slab). Implemented in the \texttt{equivalencies} array
where \texttt{equivalencies(i)} points to the canonical index. Tobias's
earlier modification to enumlib.

\item[Concentration range adjustment]
The bookkeeping that adjusts a user-supplied concentration range to
account for atoms locked into inactive sites. If a species appears
\emph{only} on inactive sites, its range is zeroed out; otherwise the
range is re-scaled.
(See \texttt{notes\_cRangeAdjustment.pdf}.)
\end{description}

\section{Misuse / scaling}

\begin{description}[leftmargin=2em, style=nextline]
\item[Pre-flight count]
Computing the number of inequivalent configurations \emph{without}
generating them, so a user can decide whether to run a full
enumeration. The Pólya enumeration theorem provides closed-form
counts via cycle indices and Burnside averaging.

\item[Pólya counting (full / unrestricted)]
Counts inequivalent colorings of $G$ over all concentrations.
Cycle-index formula: $\frac{1}{|H|} \sum_{\rho \in H}
k^{c(\rho)}$ where $c(\rho)$ is the number of cycles. Implemented as
the \texttt{polya=true} short-circuit in Fortran's
\texttt{gen\_multilattice\_derivatives}.

\item[Pólya counting at fixed concentration]
Counts inequivalent colorings with target multiplicities. Formula
(Hart-Nelson-Forcade 2012, Appendix A.2):
\[
\frac{1}{|H|} \sum_{\rho \in H} \sum_{M \in \omega(\rho)}
\prod_{i=1}^t \binom{c_i}{M_{i,1}, \ldots, M_{i,k}}
\]
where the cycle structure of $\rho$ is $(c_1, m_1), (c_2, m_2), \ldots$
and $M$ ranges over the matrices of color-multiplicities-per-cycle that
satisfy the constraints.

\item[Crossing-out algorithm (enum3)]
The 2008 + 2012 algorithm: enumerate the labeling space via a perfect
hash, mark duplicates by orbits under the symmetry group's action,
keep one representative per orbit. Memory-bound: needs $O(k^n)$ or
$O(C)$ flags. Best for site-restricted enumerations and small
multinomials.

\item[Recursive-stabilizer / tree-based algorithm (enum4)]
Morgan-Hart-Forcade 2017 algorithm. Builds an enumeration tree;
prunes via the stabilizer subgroup at each layer. Avoids the
$O(k^n)$ memory by streaming. Best for very large multinomials.

\item[\texttt{origCrossOutAlgorithm} flag]
Runtime toggle in the Fortran driver between enum3 and enum4. The
Julia rewrite should make this an explicit algorithm choice, not a
hidden boolean.
\end{description}

\end{document}
